Las técnicas se seleccionaron a partir del estado del arte (Cap.~3) y se vinculan a los factores
del Cap.~2:

\begin{itemize}
  \item \textbf{OE1:} \textit{parsing} OpenAPI~3.x; extracción de \textit{paths}, operaciones y
        esquemas; alineación con elementos BIAN en representación intermedia (Mainas et al., 2023
        como referencia de enriquecimiento semántico).
  \item \textbf{OE2:} \textit{schema matching} según taxonomía Shvaiko y Euzenat (2005): nombre,
        ruta, profundidad de esquema, operación HTTP.
  \item \textbf{OE3:} similitud semántica por embeddings y/o ontología (El Bayadh et al., 2015;
        Chandrasekaran y Mago, 2022).
  \item \textbf{OE4:} agregación ponderada $\alpha S + \beta E + \gamma C$ y clasificación por
        umbrales fijos (§1.4).
\end{itemize}

Herramientas previstas: Python para prototipo, validadores OpenAPI, bibliotecas de similitud textual;
detalle de implementación fuera de alcance de Guía~1.
